\section{Forward positive-derivative centers do not depend on initial ones}
Let $X=1^\alpha U$ be a normalized formal arm, $\alpha\ge1$,
$h=|U|\ge1$, every capacity at least two, and $W=\sum U_i$.
No ordering of the capacities is needed in this section.

\begin{theorem}[Forward derivative layer bound]
Every height layer of $D_n=(F_n-F_{n-1})^+$, $D_0=D_1=0$,
has endpoint sum at most $W+h+1$. In particular its first mode
satisfies
\[
                         \ell_D\le\left\lfloor\frac{W+h+1}{2}\right\rfloor.
\]
The bound is independent of $\alpha$. The first-mode bound is
attained on an unbounded actual family.
\end{theorem}
\begin{proof}
Define the infinite-initial-ones shifted derivative $d^\infty$
by starting at zero and replacing $d$ on append $t$ by
$[t-1]+S_td$. Direct induction on the exact arm recurrence gives
\[
 (1-z)F=1+z^2d^\infty-z^{\alpha+h+1}\prod_i[U_i],
 \qquad
 D_{n+2}=\left(d^\infty_n-[z^n]z^{\alpha+h-1}\prod_i[U_i]\right)^+.
                                                               \tag{1}
\]
The positive expansion is
\[
        d^\infty=\sum_{i=1}^h [U_i-1]\prod_{j>i}S_{U_j}.
\]
Each summand is symmetric unimodal, with endpoint sum
$\sum_{j\ge i}U_j+h-i-2\le K:=W+h-3$.
Thus for $2n\ge K$ it has coefficient at $K-n$ at least
the coefficient at $n$. Summing gives the same reflection
for $d^\infty$.

The subtracted interval product in (1) has endpoint sum
$2\alpha+W+h-2=K+2\alpha+1>K$.
Its reflection comparison goes in the opposite direction:
its coefficient at $K-n$ is at most the one at $n$.
Subtracting therefore preserves the upper-center reflection,
as does taking positive parts.
All these comparisons use zero-extended coefficients, so
support boundaries cause no exception.

The already proved forward positive-derivative unimodality
holds for every formal capacity order. For any height layer,
reflect its right endpoint $R$ about $K/2$; the reflected rank
$K-R$ belongs to the same layer, and its left endpoint is no
larger. Thus shifted endpoints sum to at most $K$.
Undoing the shift by two gives the bound $K+4=W+h+1$.
The first maximum cannot start to the right of the integer
midpoint, proving the mode estimate.

For sharpness take $U=(t,t,t)$, $t=2m+1\ge3$, and
$\alpha\ge3t$. The correction in (1) starts after the support
of $d^\infty$. Here
\[
 d^\infty=[t-1]+S_t[t-1]+S_t^2[t-1].
\]
At shifted $j=3m+1$ the last summand has increment two,
the middle one increment minus one, and the first one is zero;
the total increment is one. At $j+1$ the corresponding
increments are zero, minus one, and zero, giving minus one.
The interval-product increments follow directly from
$\Delta([t]^2[t-1])_k=(k+1)-2(k-t+1)^+-(k-t+2)^+$
at $k=3m-1,3m$. Unimodality then makes $j$ the first mode.
Its unshifted value is $3m+3=\lfloor(3t+4)/2\rfloor$,
exactly the stated bound.
\end{proof}

Combining this estimate with the complete conditional C window
theorem gives an explicit sufficient append condition independent
of initial ones:
$t+\max(4,\kappa-1)\ge\lfloor(W+h+1)/2\rfloor$.
This remains a conditional C transfer, not a general original UI proof.
